\section{导论：RAG 是系统问题，不是“搜一下再拼 Prompt”}

Douwe Kiela 的讲座从语言模型历史出发，逐层增加 retrieval、reranking、fusion、joint training、active retrieval 与 tool use。整场课的核心不是某个产品化 pipeline，而是一个研究问题：外部知识怎样进入生成过程，哪些组件应被冻结或训练，证据何时检索、如何排序、在哪里融合，以及系统怎样证明输出真正受证据约束。

\lecturefigure{slide-01-title.jpg}{Retrieval Augmented Language Models 课堂标题页}{Stanford 课堂视频 00:01:39}

\begin{knowledgebox}{读图：标题中的 Augmented 改变模型边界}
普通 language model 把主要知识压进参数；retrieval-augmented language model 在推理或训练时访问外部文档、向量索引、搜索引擎或其他工具。增强的不只是输入长度，而是可更新的 non-parametric memory 与模型之间的接口。
\end{knowledgebox}

\begin{importantbox}{本讲的统一问题}
设输入为 $x$、外部语料为 $\mathcal{D}$、检索器为 $R$、生成器为 $G$，最简流程是
\[
z_{1:k}=R(x,\mathcal{D}),\qquad
y=G(x,z_{1:k}).
\]
真正困难发生在两个等号内部：$R$ 怎样学习“对 $G$ 有用”的证据，$G$ 是否会使用 $z_{1:k}$，以及整个系统怎样在质量、成本、更新和可信度之间取舍。
\end{importantbox}

\subsection{证据边界与阅读方法}

本讲发生在 2023 年 12 月。讲者对 DRAGON、vector database、retrieval hardware 与未来产业结构的判断保留为课堂时间点意见；RAG、REALM、FiD、RETRO、Atlas、RePlug 等机制则回到原论文解释。旧稿中的 dashboard、incident runbook、gate dropout、drift monitor 和地域治理并无课堂证据，全部移除。

\subsection{本章小结}

RAG 的教学主线是从“外部文档进入 prompt”走向“检索、表示、融合与生成共同优化”。后续章节先说明为什么需要外部记忆，再用 train/test taxonomy 比较架构，最后讨论 RAG 2.0 的系统化方向。

